\documentclass{acmart}
\usepackage{mathtools}
\usepackage{mathpartir}

\usepackage{macros}

\begin{document}

\section{Language}

We work with a Python-like Language (adapted from
\cite{rak-amnouykit_python_2020}). A program is made up of declarations,
statements, and expressions.\\
\[
  d_1, d_2 \Coloneq \funDecl{f}{\vv{x}}{d_1}{d_2} \mid \varDecl{x}{e}{d_1} \mid
  s
\]

\[
  s \Coloneq x = e \mid x.attr = e \mid \ifStmt{e}{s_1}{s_2} \mid \returnStmt e
  \mid \vv{s}
\]

\[
  e_1, e_2 \Coloneq \num{n} \mid \true \mid \false \mid \dots \mid x
  \mid e.attr \mid e_1[e_2] \mid e_1(\vv{e_2}) \mid
  \listExp{e_1, \dots, e_n} \mid \dictExp{e_1: e_1', \dots, e_n:e_n' } \mid
  \newExp
\]
\\
Python's syntax does not distinguish between variable declaration, and variable
assignment, but we choose to here for simplicity. For similar reasons, we also
enforce that declarations all occur before the main body of a function or
program. Although this is not enforced by the grammar, we will assume that the
body of a function contains at most one return statement and that this will be
the last statement in the body. Any hypothetical program not following
these conventions / restrictions could be converted into a program that fits
into our grammar by a simple code transformation.

The \(\dots\) in the grammar for expressions serves as a placeholder for
a range of simple constants that are present in the language.

\section{Types}

We work with a fairly standard set of types for Python, defined by the following
grammar.
\\
\[
  A, B \Coloneq \intty \mid \boolty \mid \listty \mid
  \dictty \mid \anyty \mid \vv{A} \to B \mid [\vv{x:A}] \mid \nonety
\]
\\
The type \(\nonety\) might be expected to be related to the Python keyword of
the same name, but for our purposes it is just used as the return type for
functions that do not return a value. Our row types \([\vv{x:A}]\) are
non-standard, as they represent an upper bound on the attributes available
for an object, rather than the usual lower bound interpretation.

\begin{mathpar}
  \infer[Var]{ }{\Gamma, x:A \types x:A} \and

  \infer[Num]{ }{\Gamma \types \num{n}:\intty} \and

  \dots \and

  \infer[Dict]
  {(\forall i \in [1,n])\,\Gamma \types e_i: \anyty \\
  \Gamma \types e_i':\anyty}
  {\Gamma \types \dictExp{e_1: e_1', \dots, e_n:e_n' }:\dictty } \and

  \infer[List]
  {(\forall i \in [1,n])\,\Gamma \types e_i: \anyty}
  {\Gamma \types \listExp{e_1, \dots, e_n}:\listty } \and

  \infer[New]
  { }
  {\Gamma \types \newExp : [\vv{x:A}] } \and

  \infer[Ind-1]
  {\Gamma \types e_1: \dictty}
  {\Gamma \types e_1[e_2]: \anyty} \and

  \infer[Ind-2]
  {\Gamma \types e_1: \listty \\ \Gamma \types e_2:\intty}
  {\Gamma \types e_1[e_2]: \anyty} \and

  \infer[Attr]
  {\Gamma \types e: [attr:A, \vv{x:A}] }
  {\Gamma \types e.attr : A}

  \infer[Call]
  {\Gamma \types e: \vv{B_1, \dots B_n} \to A \\
  (\forall i \in [1,n]) \,\, \Gamma \types e_i : B_i}
  {\Gamma \types e(\vv{e_1, \dots e_n}):A}

  \\\\\\\\
  \infer[Assign-Var]
  {\Gamma, x:A \types e:A}
  {\Gamma, x:A \types x = e \rhd \nonety}

  \infer[Assign-Attr]
  {\Gamma, x:[\vv{y:B}, attr:A] \types e:A}
  {\Gamma, x:[\vv{y:B}, attr:A] \types x.attr = e \rhd \nonety}

  \dots

  \infer[If]
  {\Gamma \types e: \boolty \\ \Gamma \types s_1 \rhd A \\
  \Gamma \types s_2 \rhd B}
  {\Gamma \types \ifStmt{e}{s_1}{s_2} \rhd \nonety} \and

  \infer[Return]
  {\Gamma \types e:A}
  {\Gamma \types \returnStmt e \rhd A} \and

  \infer[Seq]
  {\Gamma \types s:B \\ \Gamma \types \vv{s}:A}
  {\Gamma \types s \, \vv{s} \,\rhd \,A}
  \\\\\\\\

  \infer[Decl]
  {\Gamma \types e:B  \\ \Gamma, x:B \types d_1 \rhd A}
  {\Gamma \types \varDecl{x}{e}{d_1} \rhd A} \and
  \infer[Fun]
  {
    \Gamma, x_1:A_i, \dots, x_n:A_n, f: \vv{A_1 \dots A_n} \to B  \types
    d_1 \rhd B \\
    \Gamma, f:\vv{A_1 \dots A_n} \to B \types d_2 \rhd A
  }
  {\Gamma \types \funDecl{f}{\vv{x_1, \dots, x_n}}{d_1}{d_2} \rhd A}
\end{mathpar}
\\\\

\bibliography{refs}{}
\bibliographystyle{plain}
\end{document}
